\documentclass[11pt]{article}
\usepackage[margin=0.8in]{geometry}
\usepackage[T1]{fontenc}
\usepackage{lmodern,amsmath,amssymb,amsthm,booktabs,array}
\usepackage[colorlinks=true,linkcolor=blue,urlcolor=blue,citecolor=blue]{hyperref}
\hypersetup{pdftitle={Cyclotomic cosets, a triangle, and reciprocal cube roots},pdfauthor={Blake Stone},pdfsubject={Six recorded source assertions in three families, with complete Lean proofs}}
\setlength{\parindent}{0pt}
\setlength{\parskip}{0.38em}
\setlength{\emergencystretch}{1em}
\newtheorem{theorem}{Theorem}[section]
\newtheorem{lemma}[theorem]{Lemma}
\newtheorem{proposition}[theorem]{Proposition}
\newtheorem{corollary}[theorem]{Corollary}
\newcommand{\N}{\mathbb N}
\newcommand{\Z}{\mathbb Z}
\newcommand{\Q}{\mathbb Q}
\newcommand{\R}{\mathbb R}
\begin{document}
\begin{center}
{\LARGE\bfseries Cyclotomic cosets, a triangle,\\[3pt]
and reciprocal cube roots}\par
\vspace{0.5em}
{\large Six source assertions in three families: complete Lean proofs}\par
\vspace{0.5em}
{\large Blake Stone}\par
Proofs developed with ChatGPT\par
November 27, 2026
\end{center}

\textbf{Abstract.}
We prove six specific conjectural or empirical assertions recorded in the
On-Line Encyclopedia of Integer Sequences, organized into three mathematical
families. First, exact prime power orders and a proved orbit-count formula
establish Barker's four quadratic counts of cyclotomic cosets modulo $10^n$,
for multipliers $3,7,9,13$. Second, a weighted Delannoy representation proves
Israel's generating function and recurrence for the subdiagonal A052392 of
Kimberling's triangle, preserving its exceptional initial entry. Third, exact
cubic inequalities prove the empirical generating function of A184534 from
its literal reciprocal fractional-root definition, and identify a displaced
index in the printed closed form. All conclusions are verified in Lean from
source definitions. Each package includes pinned sources, two separate clean
builds, and a review by an agent outside that package's authoring team.
Classical methods and known related results retain their attribution.
The bounded prior-work screen does not establish global priority.

\begin{center}
\begin{tabular}{@{}l>{\raggedright\arraybackslash}p{4.2in}@{}}
\toprule
Family & Exact scope\\
\midrule
A220018/A220019 & Quadratic orbit counts for multipliers $3$ and $7$.\\
A220020/A221855 & Quadratic orbit counts for multipliers $9$ and $13$.\\
A052392 & One subdiagonal family: empirical algebraic generating function and recurrence.\\
A184534 & One reciprocal-root family: empirical generating function, with corrected closed-form indexing.\\
\bottomrule
\end{tabular}
\end{center}

The four coset formulas are closely related specializations, counted as four
recorded assertions within one family. Equivalent recurrences, supporting
lemmas, and indexing corrections are not counted separately. General
cyclotomic-coset theory, weighted Delannoy methods, and root comparisons are
established tools, not claims of new general theories.

\section{Cyclotomic cosets modulo powers of ten}\label{sec:cosets-ten}

For $q$ coprime to $10$, let $C_q(n)$ be the number of distinct
orbits of the permutation $x\mapsto qx$ on $\mathbb Z/10^n\mathbb Z$.
Thus the orbit containing $x$ is the finite set of all $q^kx$ with
$k\ge0$. The zero orbit is included. V.~Raman's sequences
A220018, A220019, A220020, and A221855 correspond respectively to
$q=3,7,9,13$. Colin Barker recorded the following conjectures on
13--14 April 2013 \cite{oeis220018,oeis220019,oeis220020,oeis221855}.

\begin{theorem}\label{thm:cosets-ten}
For every $n\ge3$,
\[
 C_3(n)=8n^2-14n-1,\qquad
 C_9(n)=4(4n^2-7n-1),\qquad
 C_{13}(n)=2n(4n-7).
\]
For every $n\ge4$,
\[
 C_7(n)=5(16n^2-60n+37).
\]
\end{theorem}

The ingredients are standard orbit counting, the Chinese remainder
theorem, and lifting the exponent. A much more general classification
of cyclotomic cosets and their sizes appears in Zhu--Liu--Wu
\cite{zhu-liu-wu-cosets}. The present argument gives the explicit
specialization required by these four recorded formulas, with a
complete Lean proof from the literal orbit definition.

\begin{lemma}[Divisor formula]\label{lem:cosets-divisors}
For positive $m$ coprime to $q$, the number of orbits of multiplication
by $q$ on $\mathbb Z/m\mathbb Z$ is
\[
 \sum_{d\mid m}\frac{\varphi(d)}{\operatorname{ord}_d(q)}.
\]
Here $\varphi(1)=\operatorname{ord}_1(q)=1$.
\end{lemma}
\begin{proof}
The additive cyclic group $\mathbb Z/m\mathbb Z$ has exactly
$\varphi(d)$ elements of additive order $d$ for each divisor $d$ of
$m$. Multiplication by a unit preserves additive order. If $x$ has
additive order $d$, then
\[
 q^kx=x\quad\Longleftrightarrow\quad(q^k-1)x=0
 \quad\Longleftrightarrow\quad d\mid q^k-1.
\]
The least positive such $k$ is $\operatorname{ord}_d(q)$.
The $\varphi(d)$ elements in the stratum therefore partition into
orbits of that common size, giving the displayed quotient.
Summing over all divisors counts every orbit exactly once.
The stratum $d=1$ consists only of zero.
\end{proof}

Write $t_+=\max(t,0)$. The exact orders needed in the proof are
\begin{equation}\label{eq:cosets-orders-two}
\begin{array}{c|c|c}
q& i=0,1&i\ge2\\\hline
3&1&2\cdot2^{(i-3)_+}\\
7&1&2\cdot2^{(i-4)_+}\\
9&1&2^{(i-3)_+}\\
13&1&2^{(i-2)_+}
\end{array}
\end{equation}
and, for $j\ge1$,
\begin{equation}\label{eq:cosets-orders-five}
\begin{array}{c|c}
q&\operatorname{ord}_{5^j}(q)\\\hline
3&4\cdot5^{j-1}\\
7&4\cdot5^{(j-2)_+}\\
9&2\cdot5^{j-1}\\
13&4\cdot5^{j-1}.
\end{array}
\end{equation}
The orders at $j=0$ are $1$.

\begin{proof}[Proof of the order formulas]
The orders of $3,7$ modulo $4$ are $2$, and their exact valuations are
$v_2(3^2-1)=3$ and $v_2(7^2-1)=4$. Lifting the exponent gives
\[
 v_2(3^{2k}-1)=3+v_2(k),\qquad
 v_2(7^{2k}-1)=4+v_2(k).
\]
For $i\ge2$, an exponent producing $1$ modulo $2^i$ must be even.
The displayed laws therefore show that the least possible exponent is
$2\cdot2^{(i-3)_+}$ or $2\cdot2^{(i-4)_+}$, respectively.
For $9$ and $13$, which are $1$ modulo $4$, the analogous laws are
\[
 v_2(9^k-1)=3+v_2(k),\qquad
 v_2(13^k-1)=2+v_2(k).
\]
They give the other two rows, with the small moduli verified directly.

Modulo $5$, the orders of $3,7,9,13$ are $4,4,2,4$, respectively.
The exact valuations of $q^e-1$ at these exponents are $1,2,1,1$.
In particular $7^4-1=2400=5^2\cdot96$ is divisible by $25$ but
not $125$. For each of the four pairs $(q,e)$, lifting the exponent
gives
\[
 v_5(q^{ek}-1)=v_5(q^e-1)+v_5(k).
\]
Any exponent producing $1$ modulo $5^j$ is a multiple of $e$;
the smallest possible multiplier is the power of $5$ displayed in
\eqref{eq:cosets-orders-five}. These statements apply to every
positive exponent, not only a finite list of moduli.
\end{proof}

Every divisor of $10^n$ has a unique representation $2^i5^j$ with
$0\le i,j\le n$. By the Chinese remainder theorem,
\[
 \operatorname{ord}_{2^i5^j}(q)
 =\operatorname{lcm}\bigl(\operatorname{ord}_{2^i}(q),
                         \operatorname{ord}_{5^j}(q)\bigr).
\]
Put
\[
 T_q(i,j)=\frac{\varphi(2^i5^j)}{\operatorname{ord}_{2^i5^j}(q)}.
\]
Then Lemma~\ref{lem:cosets-divisors} becomes
\begin{equation}\label{eq:cosets-grid}
 C_q(n)=\sum_{i=0}^n\sum_{j=0}^n T_q(i,j).
\end{equation}
Multiplicativity of $\varphi$, together with
$\varphi(2^{i+1})=2^i$ and $\varphi(5^{j+1})=4\cdot5^j$,
gives the following complete table:
\[
\begin{array}{c|rrrrrr}
 &i=0&i=1&i=2&i=3&i=4&i\ge5\\\hline
T_3(i,0)&1&1&1&2&2&2\\
T_3(i,j),\ j\ge1&1&1&2&4&8&8\\
T_7(i,0)&1&1&1&2&4&4\\
T_7(i,1)&1&1&2&4&8&16\\
T_7(i,j),\ j\ge2&5&5&10&20&40&80\\
T_9(i,0)&1&1&2&4&4&4\\
T_9(i,j),\ j\ge1&2&2&4&8&16&16\\
T_{13}(i,0)&1&1&2&2&2&2\\
T_{13}(i,j),\ j\ge1&1&1&2&4&8&8
\end{array}
\]
For example, if $i=r+4$ and $j=s+1$, the orders for base $3$ are
$4\cdot2^r$ and $4\cdot5^s$. Since $2^r$ and $5^s$ are coprime,
their least common multiple is $4\cdot2^r5^s$. Dividing
$2^{r+3}\cdot4\cdot5^s$ by it gives $8$. The other infinite
tails follow by the same exact cancellations.

\begin{proof}[Proof of Theorem~\ref{thm:cosets-ten}]
The sum through index $n$ of $1,1,2,4,8,8,\ldots$ is
$8n-16$ for $n\ge3$. The sum through index $n$ of
$1,1,2,4,8,16,16,\ldots$ is $16n-48$ for $n\ge4$.
The four $j=0$ rows sum respectively to $2n-1$, $4n-7$,
$4n-4$, and $2n$, in the ranges required by the theorem.
Consequently \eqref{eq:cosets-grid} gives
\begin{align*}
 C_3(n)&=(2n-1)+n(8n-16)=8n^2-14n-1,\\
 C_7(n)&=(4n-7)+(1+5(n-1))(16n-48)
        =80n^2-300n+185,\\
 C_9(n)&=(4n-4)+2n(8n-16)=16n^2-28n-4,\\
 C_{13}(n)&=2n+n(8n-16)=8n^2-14n.
\end{align*}
These are exactly the four conjectured polynomials.
\end{proof}

\paragraph{Formal scope and attribution.}
The Lean definition counts the quotient of $\mathbb Z/10^n\mathbb Z$
by the action of the subgroup generated by the unit $q$.
The theorem \path{mem_coset_iff} proves that membership in these
orbits is equivalent to $y=q^kx$ for some natural $k$.
The divisor formula is then proved using orbit sizes and additive
orders. Subsequent modules prove the divisor-to-exponent bijection,
all prime power orders by exact valuations, every entry and tail of
the table, and all four sums. Thus the polynomial is the conclusion,
not the definition of any formal source sequence.

The four endpoint theorems in \path{Proof.lean} use exactly the
source ranges $n>2$ or $n>3$ and have only the standard Lean logical
axioms as dependencies. All six modules are explicit build roots.
The source entries and earlier generators already contain the
divisor algorithm or polynomial patterns; the general methods are
classical and are not counted as new results. This package contains
four related source-assertion proofs in one family. The separately
recorded bounded prior-work search found no earlier explicit
settlement of these four formulas and does not establish worldwide
priority.


\section{A052392: a subdiagonal of Kimberling's triangle}
\label{sec:a052392}

\subsection{Source statement and indexing}

OEIS A052392 \cite{oeis052392} defines $a_n=T(2n+1,n)$, starting with $a_0=1$.
Robert Israel's November 26, 2018 contribution records an empirical
generating function and an empirical order-four recurrence. We prove both
as one source family. The algorithm attached to that entry defines
$T(r,0)=T(r,r)=1$, $T(r,k)=0$ for $k>r$, and $T(2,1)=3$; its remaining
entries satisfy
\begin{equation}\label{eq:a052392-triangle}
 T(r,k)=T(r-1,k-1)+2T(r-2,k-1)+T(r-1,k)
 \qquad(r\ge3,\ 0<k<r).
\end{equation}
The exceptional value $T(2,1)=3$ is essential. Applying the interior rule
there would instead produce $4$ and a different sequence.

The associated entry A054120 originally presents this triangle through
strip sums. Russell Jay Hendel's June 2, 2015 comment in that entry already
proves \eqref{eq:a052392-triangle} from the strip-sum definition and
explains its exceptional initial entry. That equivalence is credited
prior work. Our formal definition follows the literal algorithm attached
to A052392, rather than defining the sequence by the formula to be proved.

\begin{theorem}\label{thm:a052392-main}
Let $a_n=T(2n+1,n)$ with the preceding boundary conditions and recurrence.
In the formal power-series ring $\mathbb Q[[x]]$,
\begin{equation}\label{eq:a052392-gf}
 A(x):=\sum_{n\ge0}a_nx^n
 =\frac{2x^2-3x+1}{2x\sqrt{1-8x+4x^2}}+\frac{x-1}{2x},
\end{equation}
where the square root has constant coefficient $1$. The numerator of
the combined expression vanishes at zero, so the displayed expression
is a power series. Furthermore, for every $n\ge0$,
\begin{equation}\label{eq:a052392-rec}
\begin{split}
 8na_n-(32+28n)a_{n+1}+(78+30n)a_{n+2}\\
 {}-(45+11n)a_{n+3}+(5+n)a_{n+4}=0.
\end{split}
\end{equation}
\end{theorem}

\subsection{An independent array representation}

Define the weighted Delannoy array by
\[
 D(0,n)=D(m,0)=1,\qquad
 D(m+1,n+1)=D(m,n+1)+D(m+1,n)+2D(m,n).
\]
This definition determines every entry by induction on $m+n$ and gives
$D(m,n)=D(n,m)$. Write $S(m,n)=T(m+n,n)$. Thus $S$ has the same boundary
values and the same interior recurrence, except that $S(1,1)=3$.

\begin{lemma}\label{lem:a052392-difference}
For all $m,n\ge0$,
\[
 S(m+1,n+1)=D(m+1,n+1)-D(m,n).
\]
\end{lemma}
\begin{proof}
At $(m,n)=(0,0)$ the right side is $4-1=3$. Along the first interior
row, $D(1,n)=3n+1$, whereas $S(1,n)=3n$ for $n\ge1$; these identities
follow immediately by induction from the boundary rule. Symmetry gives
the first interior column. At all remaining positions, subtracting the
recurrences for $D(m+1,n+1)$ and $D(m,n)$ gives exactly the recurrence
for $S(m+1,n+1)$. Induction on $m+n$ now proves the result.
\end{proof}

Let
\[
 d_n=D(n,n),\qquad b_n=D(n+1,n).
\]
Then $d_0=b_0=a_0=1$, $d_1=4$, and Lemma
\ref{lem:a052392-difference} gives
\begin{equation}\label{eq:a052392-a-b}
 a_{n+1}=b_{n+1}-b_n\qquad(n\ge0).
\end{equation}
By symmetry and the array recurrence,
\begin{equation}\label{eq:a052392-d-b}
 d_{n+1}=2b_n+2d_n\qquad(n\ge0).
\end{equation}

\subsection{The diagonal recurrence}

\begin{lemma}\label{lem:a052392-balance}
For every $m,n\ge0$,
\begin{equation}\label{eq:a052392-balance}
 (m+1)D(m+1,n)-(n+1)D(m,n+1)=2(n-m)D(m,n).
\end{equation}
\end{lemma}
\begin{proof}
Let $E(m,n)$ be the left side minus the right side of
\eqref{eq:a052392-balance}. On the axes, the formula
$D(1,n)=3n+1$ and symmetry give $E(0,n)=E(m,0)=0$.
Expanding three instances of the defining recurrence for $D$ gives
\[
 E(m+1,n+1)=E(m,n+1)+E(m+1,n)+2E(m,n).
\]
For an explicit check, subtract the right side from the left; the
remaining expression is
\[
 2(m-n)\bigl(D(m+1,n+1)-D(m,n+1)-D(m+1,n)-2D(m,n)\bigr)=0.
\]
Induction on $m+n$ therefore gives $E(m,n)=0$ everywhere.
\end{proof}

Apply \eqref{eq:a052392-balance} at $(m,n)=(n+1,n)$, and use
\[
 b_{n+1}=D(n+1,n+1)+D(n+2,n)+2D(n+1,n).
\]
The result is
\begin{equation}\label{eq:a052392-brec}
 (n+2)b_{n+1}=(2n+3)d_{n+1}+2(n+1)b_n.
\end{equation}
Eliminating $b_n,b_{n+1}$ with \eqref{eq:a052392-d-b} yields
\begin{equation}\label{eq:a052392-drec}
 (n+2)d_{n+2}=4(2n+3)d_{n+1}-4(n+1)d_n\qquad(n\ge0).
\end{equation}
Likewise, \eqref{eq:a052392-a-b} and \eqref{eq:a052392-d-b} give
\begin{equation}\label{eq:a052392-source-central}
 2a_{n+1}=d_{n+2}-3d_{n+1}+2d_n\qquad(n\ge0).
\end{equation}
Every relation so far is an identity for the independently defined arrays.

\subsection{The generating function}

Put $D(x)=\sum_{n\ge0}d_nx^n$ and $Q(x)=1-8x+4x^2$.
Equation \eqref{eq:a052392-drec}, together with $d_0=1$ and $d_1=4$,
is equivalent coefficient by coefficient to
\begin{equation}\label{eq:a052392-ode}
 QD'=(4-4x)D.
\end{equation}
Indeed, the coefficient of $x^j$ is
$(j+1)d_{j+1}=(8j+4)d_j-4jd_{j-1}$ for $j\ge1$,
and the constant coefficient says $d_1=4d_0$.
Formal differentiation now gives
\[
 (QD^2)'=(-8+8x)D^2+2D\,QD'=0.
\]
Over $\mathbb Q[[x]]$, a series with zero derivative is constant, and the
constant coefficient here is $1$. Thus
\begin{equation}\label{eq:a052392-square}
 QD^2=1,\qquad D(0)=1.
\end{equation}
These conditions characterize $D$ uniquely: if $QE^2=1$ and $E(0)=1$,
then $Q\ne0$ gives $(E-D)(E+D)=0$. The ring is a domain and $E+D$ has
constant coefficient $2$, so $E=D$. Consequently $D=1/\sqrt Q$ with
the normalized branch. This argument proves existence and uniqueness;
it does not assume a square-root coefficient expansion.

Multiply \eqref{eq:a052392-source-central} by $x^{n+2}$ and sum over
$n\ge0$, using $a_0=d_0=1$ and $d_1=4$. One obtains
\begin{equation}\label{eq:a052392-cleared}
 2xA(x)=(1-3x+2x^2)D(x)-1+x
       =(1-x)\bigl((1-2x)D(x)-1\bigr).
\end{equation}
Substitution of $D=1/\sqrt Q$ gives exactly
\eqref{eq:a052392-gf}. In particular the right side of
\eqref{eq:a052392-cleared} has zero constant coefficient. Multiplication
by $2x$ is injective in $\mathbb Q[[x]]$, so the cleared identity loses
no information about $A$.

\subsection{The exact source recurrence}

For clarity, an explicit algebraic certificate gives the recurrence
without relying on a guessed differential equation for $A$. Define
\[
 R_j=(j+2)d_{j+2}-4(2j+3)d_{j+1}+4(j+1)d_j.
\]
By \eqref{eq:a052392-drec}, every $R_j=0$.
Let $L_n$ denote the left side of \eqref{eq:a052392-rec}.
Substitute \eqref{eq:a052392-source-central} into the five terms of
$2L_{j+1}$. Collecting coefficients of $d_j,\ldots,d_{j+6}$ gives the
polynomial identity
\[
 2L_{j+1}=4R_j-12R_{j+1}+13R_{j+2}-6R_{j+3}+R_{j+4}=0
 \qquad(j\ge0).
\]
Hence $L_n=0$ for every $n\ge1$. Finally the source array gives
$(a_0,a_1,a_2,a_3,a_4)=(1,6,39,261,1779)$, and
\[
 L_0=-32\cdot6+78\cdot39-45\cdot261+5\cdot1779=0.
\]
This proves \eqref{eq:a052392-rec} at every source index and completes
the theorem.

\subsection{Formal coverage and provenance}

The Lean development defines the source array and the weighted Delannoy
array independently. It proves their relation, the balance identity,
the diagonal recurrence, the normalized formal reciprocal-square-root
identity and uniqueness, the cleared generating function, and the exact
order-four recurrence for every natural index. The single audited
endpoint is \texttt{A052392.full\_resolution}. Its dependencies are only
the standard Lean axioms \texttt{propext}, \texttt{Classical.choice}, and
\texttt{Quot.sound}. All four local modules are explicit default build
roots; the compiler and mathlib revision are pinned.

The frozen sources are OEIS A052392 revision 12 (April 24, 2025) and
A054120 revision 27 (January 26, 2022), retrieved from immutable
OEIS-data commits. Weighted Delannoy methods and generalized Pascal
triangles are classical; see Paul Barry, \emph{On Integer-Sequence-Based
Constructions of Generalized Pascal Triangles}, Journal of Integer
Sequences \textbf{9} (2006), Article 06.2.4. The exact and semantic
prior-work check is documented separately, including known source
implementations and the current ProveIt and trureturing repositories.
No proof of the precise A052392 empirical assertions was found within
that bounded search. We assert resolution of this recorded source
family, without asserting worldwide priority for its underlying methods.


\section{A reciprocal fractional cube root and its generating function}

Kimberling's sequence A184534 is defined, for $n\ge1$, by
\[
 a(n)=\left\lfloor\frac1{\{(n^3+4)^{1/3}\}}\right\rfloor,
\]
where braces denote fractional part. Barker recorded an empirical rational
generating function and a separate empirical closed form in 2012
\cite{oeis184534}. We prove the generating function and identify the
one-index displacement in the printed closed form.

\begin{theorem}
The source sequence satisfies
\[
 a(1)=1,\qquad a(2)=3,\qquad a(3)=7,\qquad
 a(n)=\left\lfloor\frac{3n^2}{4}\right\rfloor\quad(n\ge4).
\]
Its ordinary generating function is exactly
\begin{equation}\label{a184534:gf}
 \sum_{n\ge1}a(n)x^n=
 \frac{x(x^6-2x^5+x^4-x^2-x-1)}{(x-1)^3(x+1)}.
\end{equation}
For every $n\ge8$,
\[
 a(n)=2a(n-1)-2a(n-3)+a(n-4).
\]
\end{theorem}

\begin{lemma}\label{a184534:general}
For positive integers $h\le n$,
\[
 \left\lfloor\frac1{\{(n^3+h)^{1/3}\}}\right\rfloor
 =\left\lfloor\frac{3n^2}{h}\right\rfloor.
\]
\end{lemma}
\begin{proof}
Put $q=\lfloor3n^2/h\rfloor$. Then $q\ge3n>0$,
$hq\le3n^2$, and, by integrality, $h(q+1)\ge3n^2+1$.
First,
\[
 (nq+1)^3-(n^3+h)q^3
 =q^2(3n^2-hq)+3nq+1>0.
\]
Second,
\begin{align*}
 &(n^3+h)(q+1)^3-(n(q+1)+1)^3\\
 &\quad=(q+1)^2\bigl(h(q+1)-3n^2\bigr)-3n(q+1)-1\\
 &\quad\ge(q-3n)(q+1)+q>0.
\end{align*}
Strict monotonicity of cubing therefore gives
\[
 n+\frac1{q+1}<(n^3+h)^{1/3}<n+\frac1q.
\]
In particular the root has integer part $n$. Its positive fractional
part lies strictly between $1/(q+1)$ and $1/q$, so its reciprocal lies
strictly between $q$ and $q+1$. The floor is $q$.
\end{proof}

\begin{proof}[Proof of the theorem]
Apply Lemma~\ref{a184534:general} with $h=4$ for every $n\ge4$.
The three earlier values follow from the same two cubic comparisons,
using the following exact integer certificates:
\[
\begin{array}{c|c|c|c}
n&q&(n^3+4)q^3<(nq+1)^3&
(n(q+1)+1)^3<(n^3+4)(q+1)^3\\\hline
1&1&5<8&27<40\\
2&3&324<343&729<768\\
3&7&10633<10648&15625<15872
\end{array}
\]
These comparisons also show the root has integer part $n$ in each
case; no numerical approximation to a root is required.

For every natural $n$, the identity
$\lfloor3n^2/4\rfloor=3\lfloor n^2/4\rfloor$ follows from
$n^2\equiv0$ or $1\pmod4$. The only positive-index corrections to this
formula are one additional unit at $n=1$ and at $n=3$. Splitting into
even and odd indices gives
\[
 \sum_{n\ge0}\left\lfloor\frac{n^2}4\right\rfloor x^n
 =\frac{x^2}{(1-x)^3(1+x)}.
\]
Consequently
\[
 \sum_{n\ge1}a(n)x^n
 =x+x^3+\frac{3x^2}{(1-x)^3(1+x)}.
\]
Combining the terms gives~\eqref{a184534:gf}. The normalized denominator
is $1-2x+2x^3-x^4$ and its numerator has degree seven. Coefficient
comparison proves the recurrence from newest index eight onward.
\end{proof}

\paragraph{The source's closed-form indexing.}
The expression printed as an empirical value of $a(n)$ is
\[
 \frac{3(1-(-1)^n+4n+2n^2)}8\qquad(n>2).
\]
It equals $12$ at $n=3$, while $a(3)=7$. The expression is instead
$a(n+1)$ for every $n\ge3$. The corrected same-index expression is
\[
 a(n)=\frac{3(2n^2-1+(-1)^n)}8\qquad(n\ge4).
\]
The recurrence cutoff also matters: at newest index seven its right
side is $37$, whereas $a(7)=36$.

\paragraph{Formal conclusion and attribution.}
Four Lean modules prove the literal real-root and fractional-part
statement, the initial cubic certificates, the all-index sequence
formula, the recurrence, and the exact formal generating function in
$\mathbb Z[[X]]$. The final theorem \texttt{A184534.full\_resolution}
also identifies every positive source coefficient and the zero
constant coefficient. The source's denominator has constant term
$-1$, so the formal quotient is unambiguous. A separate Lean theorem
certifies the printed closed-form discrepancy at $n=3$.

The eventual expression already occurs in sequence generators, and
the root-comparison method is elementary. Those expressions and
methods are not claimed as new discoveries. The contribution counted
here is one complete proof of the recorded empirical generating-function
assertion, with the indexing issue documented. The general array lemma
and the recurrence are parts of the same package. The separate bounded
prior-work record found no earlier proof connecting the literal source
to this identity and does not establish worldwide priority.

\section{Verification, review, and source scope}

The three independently buildable projects contain fourteen local proof
modules. Every module is an explicit default Lake root. All projects pin
Lean 4.19.0 and mathlib commit
\begin{center}
\texttt{c44e0c8ee63ca166450922a373c7409c5d26b00b}.
\end{center}

\begin{center}
\begin{tabular}{@{}lcp{3.8in}@{}}
\toprule
Package & Modules & Formal endpoint\\
\midrule
A220018 & 6 & Four polynomial counts of literal multiplication-orbit quotients, at every stated index.\\
A052392 & 4 & The complete normalized algebraic generating function and order-four recurrence of the source triangle's subdiagonal.\\
A184534 & 4 & The literal real-root floor formula, exact source generating function, and unbounded recurrence.\\
\bottomrule
\end{tabular}
\end{center}

Every project has an author build and a separate reviewer build in a fresh
isolated directory containing only its proof sources and pinned configuration.
No local compiled proof objects are copied; only pinned external dependency
caches are shared. The logs record exact source hashes and printed axiom
audits. The principal endpoints depend only on \path{propext},
\path{Classical.choice}, and \path{Quot.sound}. No proof uses an admitted
lemma, a custom axiom, or \path{native_decide}.

A separate reviewer checks each package's mathematical argument, source
interpretation, exceptional indices, and formal theorem scope. These are
independent reviews within this collaboration, not external journal peer
review. Supplemental finite computations test transcription and implementation;
the unbounded conclusions are established by the proofs above and the Lean
kernel, not by extrapolation from sampled values.

\paragraph{Attribution and novelty scope.}
The coset count is an explicit specialization of classical order and orbit
methods; Zhu--Liu--Wu provide a much more general classification. Hendel's
source comment already establishes the triangle recurrence, and weighted
Delannoy theory is classical. The empirical formulas in the source entries
and numerical generators also predate these proofs. The contribution claimed
here is the full proof and formal verification of these exact recorded
assertions, with source indexing issues made explicit.

The source dossiers preserve immutable OEIS exports, exact and semantic
searches across both web search systems, global GitHub searches, and current
checks of the ProveIt and trureturing proof collections. Related papers and
programs were inspected for their actual theorem scope. These bounded checks
found no earlier explicit settlement of the six assertions addressed here;
search coverage and publication indexing remain incomplete. No claim of
exhaustive historical priority or external acceptance follows from this
publication. The repository publication does not edit OEIS.

\begin{thebibliography}{12}
\raggedright
\bibitem{oeis220018} V.~Raman and C.~Barker, OEIS A220018, \url{https://oeis.org/A220018}, source revision 36, 3 July 2023.
\bibitem{oeis220019} V.~Raman and C.~Barker, OEIS A220019, \url{https://oeis.org/A220019}, source revision 27, 3 July 2023.
\bibitem{oeis220020} V.~Raman and C.~Barker, OEIS A220020, \url{https://oeis.org/A220020}, source revision 26, 12 July 2025.
\bibitem{oeis221855} V.~Raman and C.~Barker, OEIS A221855, \url{https://oeis.org/A221855}, source revision 23, 3 July 2023.
\bibitem{zhu-liu-wu-cosets} L.~Zhu, J.~Liu and H.~Wu, Explicit representatives and sizes of cyclotomic cosets and their application to cyclic codes over finite fields, \emph{Finite Fields and Their Applications} \textbf{111} (2026), 102761. \url{https://arxiv.org/abs/2410.12122}; \url{https://doi.org/10.1016/j.ffa.2025.102761}.
\bibitem{oeis052392} C.~Kimberling; empirical formulas by R.~Israel, November 26, 2018.
OEIS A052392, revision 12, April 24, 2025. \url{https://oeis.org/A052392}.
\bibitem{oeis054120} C.~Kimberling; recurrence proof by R.~J.~Hendel, June 2, 2015.
OEIS A054120, revision 27, January 26, 2022. \url{https://oeis.org/A054120}.
\bibitem{barry2006} P.~Barry. \emph{On Integer-Sequence-Based Constructions of Generalized Pascal Triangles}.
Journal of Integer Sequences \textbf{9} (2006), Article 06.2.4.
\url{https://www.maths.tcd.ie/EMIS/journals/JIS/VOL9/Barry/barry91.pdf}.
\bibitem{oeis184534}
C. Kimberling; empirical formulas by C. Barker.
OEIS A184534, revision 15, May 15, 2017.
\url{https://oeis.org/A184534}.

\end{thebibliography}
\end{document}
